% Credit: pgfplots manual
\begin{frame}
    \nointerlineskip
    \begin{tikzpicture}[/pline-vec/main orientation = {90}{-60}{\t}]
        \useasboundingbox
            (-\textwidth/2,-\textheight/2)
            rectangle
            (\textwidth/2,\textheight/2);
        \pvappendsurface[
            u min = 0
            ,u max = 180
            ,u samples = 41
            ,v min = 0
            ,v max = 360
            ,v samples = 25
        ]{
            -2/15 * math.cosd(u) * (
                3 * math.cosd(v) - 30 * math.sind(u) +
                90 * math.cosd(u)^4 * math.sind(u) -
                60 * math.cosd(u)^6 * math.sind(u) +
                5 * math.cosd(u) * math.cosd(v) * math.sind(u)
            )
        }
        {
            -1/15 * math.sind(u) * (
                3 * math.cosd(v) -
                3 * math.cosd(u)^2 * math.cosd(v) -
                48 * math.cosd(u)^4 * math.cosd(v) +
                48 * math.cosd(u)^6 * math.cosd(v) -
                60 * math.sind(u) +
                5 * math.cosd(u) * math.cosd(v) * math.sind(u) -
                5 * math.cosd(u)^3 * math.cosd(v) * math.sind(u) -
                80 * math.cosd(u)^5 * math.cosd(v) * math.sind(u) +
                80 * math.cosd(u)^7 * math.cosd(v) * math.sind(u)
            )
        }
        {2/15 * (3 + 5 * math.cosd(u) * math.sind(u)) * math.sind(v)}
        \pvrendersegments
    \end{tikzpicture}
\end{frame}